Write $\mathfrak{p}=(\mathfrak{a}:x)$ with $x\notin\mathfrak{a}$. If $yz\in\mathfrak{p}$ and $z\notin\mathfrak{p}$, then $zx\notin\mathfrak{a}$ and $\mathfrak{p}\subseteq(\mathfrak{a}:zx)$. Maximality gives equality, and $y\in(\mathfrak{a}:zx)=\mathfrak{p}$. Thus $\mathfrak{p}$ is prime. Since $\mathfrak{p}=r(\mathfrak{a}:x)$, the first uniqueness theorem shows that it belongs to $\mathfrak{a}$.
